\section{小容量渐近的完整推导}
\label{app:small-capacity-proof}
\begin{proof}[命题\ref{prop:small-capacity}的证明]
令 $k=K/h$、$u=(s_B-h)/h$、$v=(\bar s_B-h)/h$，并在本附录记
$f(v)=v-\ln(1+v)$。安全边界给出 $\bar i_B/h=k-f(v)$。
完全跟踪不变量与 $\Theta$ 匹配分别等价于
\begin{equation}
 \beta\ln\frac{1+u}{1+v}=f(v),\qquad
 \frac{u}{k(\beta+u)}=
 \frac{v}{[k-f(v)](\beta+v)}.
 \label{eq:small-K-endpoint-match}
\end{equation}
因此可用 $v$ 参数化容量切换点：
\begin{equation}
 u=(1+v)\exp[f(v)/\beta]-1,\qquad
 k=\frac{u(\beta+v)f(v)}{\beta(u-v)}.
 \label{eq:small-K-parametrization}
\end{equation}
第一式在 $v=0$ 附近解析，且
\begin{align}
 u&=v+\frac{v^2}{2\beta}+\frac{v^3}{6\beta}+O(v^4),\notag\\
 k&=\beta v+\left(\frac32-\beta\right)v^2
   +\left(\beta-\frac{19}{12}+\frac1{2\beta}\right)v^3+O(v^4).
 \label{eq:small-K-param-series}
\end{align}
这里 $f(v)=v^2/2+O(v^3)$，$u-v=v^2/(2\beta)+O(v^3)$，
故第二式的奇点可去，其解析延拓满足 $k(0)=0,k'(0)=\beta>0$。
解析逆函数定理给出小正 $k$ 的唯一局部正支 $v(k)$，
且 $k-f(v)>0$、$u>v>0$；结合既有容量切换点的唯一性，这就是 $s_B$ 对应的分支。
反演得
\begin{equation}
 v=\frac k\beta-\frac{(3/2-\beta)k^2}{\beta^3}+O(k^3),\qquad
 u=\frac k\beta-\frac{(1-\beta)k^2}{\beta^3}+O(k^3),
 \label{eq:small-K-inverse}
\end{equation}
还原 $s_B-h=hu$ 即得式\eqref{eq:small-K-switch}。

由容量弧与完全跟踪弧的成本公式，
\begin{align}
 hC_B(s_B,h)
 &=\frac\beta k\int_0^u\frac{t}{(1+t)(\beta+t)}\dd t
 =\frac v{2\beta}+\frac{(2\beta-7)v^2}{12\beta^2}+O(v^3),\notag\\
 hW(s_B,K)
 &=-\ln\left(1-\frac{f(v)}k\right)
 =\frac v{2\beta}+\frac{(4\beta-15)v^2}{24\beta^2}+O(v^3).
 \label{eq:small-K-cost-series}
\end{align}
上述函数在零点均有解析延拓，故相减得到
$h\Delta J=v^2/(24\beta^2)+O(v^3)=k^2/(24\beta^4)+O(k^3)$，
即式\eqref{eq:small-K-cost}。
再多展开一阶有
\begin{equation}
 \Delta J=\frac{K^2}{24\beta^4h^3}
 \left[1-\frac{2(2-\beta)K}{\beta^2h}
       +O\bigl((K/h)^2\bigr)\right].
 \label{eq:small-K-precision}
\end{equation}
此式的余项常数仍依赖固定的 $\beta$；它说明小概率时的精度限制，
不能用来声称关于 $\beta\downarrow0$ 的一致展开。

公共容量段严格抵消，且 $c=\gamma/(\beta h)$，所以
\begin{equation}
 \gamma(T_{+,\rm fill}-T_+^*)
 =\frac\beta k\ln\left(1+\frac u\beta\right)
   +\ln\left(1-\frac{f(v)}k\right)
 =\frac1\beta+O(k).
 \label{eq:small-K-time-series}
\end{equation}
由此得式\eqref{eq:small-K-time}。
以上各函数均在固定参数下解析，其泰勒余项在足够小的邻域有界。
公共尾段差不依赖 $s_1$，故这些余项对 $s_1\ge s_B$ 的入口统一。

最后，$s_B\to h<s_1^\circ$ 保证初值族最终满足 $s_1(K)>s_B$。
利用
\[
 \mathcal C(S)=\beta\int_h^S\frac{s-h}{s(s-p)}\dd s>0\quad(S>h),
 \qquad J_{\rm fill}=\frac{\mathcal C(s_1(K))}K,
\]
及 $\mathcal C(s_1(K))\to\mathcal C(s_1^\circ)>0$，得到式\eqref{eq:small-K-relative}。
在仅有 $s_1(K)\to s_1^\circ$ 的条件下，相对节省的等价式只承诺 $1+o(1)$ 的相对余项。
又 $T_{+,\rm fill}=(cK)^{-1}\ln[(s_1(K)-p)/\ell]$ 为 $K^{-1}$ 阶，
故 $T_+^*/T_{+,\rm fill}\to1$。命题得证。
\end{proof}
